\clearpage\thispagestyle{empty}
\pdfbookmark[0]{Quellen zum Ceres-Bericht vom Juni 1802}{ceres180206-sources}
\begin{center}\textbf{Quellen zum Ceres-Bericht vom Juni 1802}\end{center}
\noindent Diese KI-Übertragung in heutiges Deutsch folgt dem Auszug aus Franz Xaver von Zachs Bericht über Ceres in \href{https://gdz.sub.uni-goettingen.de/id/PPN236007467}{Gauß’ \textit{Werken}, Band~VI} (Göttingen, 1874), S.~224. Von Zach berichtet über Gauß’ Vergleich mit den Beobachtungen und über seine Berechnung der Opposition.
\par\smallskip\noindent Der \href{https://api.digitale-sammlungen.de/iiif/image/v2/bsb10538597_00592/full/1450,/0/default.jpg}{Erstdruck vom Juni 1802} steht als Beitrag LIX in der \textit{Monatlichen Correspondenz}, Band~V, S.~576--590. Der Auszug in den \textit{Werken} reicht vom Absatz über Gauß’ Vergleich und der \href{https://api.digitale-sammlungen.de/iiif/image/v2/bsb10538597_00598/full/1450,/0/default.jpg}{Tabelle mit zwölf Zeilen auf S.~582} bis zum \href{https://api.digitale-sammlungen.de/iiif/image/v2/bsb10538597_00599/full/1450,/0/default.jpg}{Absatz über die Opposition auf S.~583}.
\par\smallskip\noindent Die Tabelle der \textit{Werke} enthält eine Spalte für die mittlere Seeberger Zeit und bezeichnet die Unterschiede mit R—B. Für den 1. und 2. April stehen dort die berechneten Deklinationen \(18\ 3\ 20.6\) und \(18\ 5\ 6.2\), im Zeitschriftendruck dagegen \(18\ 3\ 20{,}0\) und \(18\ 5\ 0{,}2\). Für den 5. April lautet die berechnete Rektaszension \(179\ 40\ 15.6\) in den \textit{Werken} und \(179\ 40\ 15{,}0\) in der Zeitschrift. Die Übertragung folgt den \textit{Werken}.
\par\smallskip\noindent Übertragung und Satz wurden mit KI erstellt. Die Digitalisate der Zeitschrift stammen von der Bayerischen Staatsbibliothek.
\clearpage
